\chapter{Résultats et Validation}

\section{Résultats du système d'inspection}
\subsection{Performance du sous-système mécanique}
\todo{Tableaux de mesures : précision angulaire, répétabilité, temps d'indexation par veine.}

\subsection{Performance du sous-système vision}
\todo{Résolution spatiale effective, qualité d'image sur M00, sensibilité détection 0.8 mm.}

\subsection{Performance du modèle IA}
\todo{Métriques finales : précision, rappel, F1-score, AUC. Matrice de confusion. Exemples d'images correctement et incorrectement classifiées.}

\subsection{Performance globale — temps de cycle}
\todo{Mesure du temps de cycle complet par pièce M00 (40 veines, 80 surfaces). Comparaison avec l'objectif $\leq 3$~min.}

\section{Validation des critères de succès}
\begin{table}[H]
\centering
\caption{Validation des critères de succès VisionGuard}
\begin{tabular}{lll}
\toprule
\textbf{Critère} & \textbf{Objectif} & \textbf{Résultat} \\
\midrule
Détection défauts $\geq 0.8$~mm  & 100\% détection  & \todo{à mesurer} \\
Temps de cycle                    & $\leq 3$~min/pièce & \todo{à mesurer} \\
Précision modèle IA               & $> 90\%$         & \todo{à mesurer} \\
Rappel modèle IA                  & $> 90\%$         & \todo{à mesurer} \\
Reproductibilité                  & Résultats stables & \todo{à évaluer} \\
\bottomrule
\end{tabular}
\end{table}

\section{Discussion}
\todo{Analyser les résultats obtenus face aux objectifs. Identifier les limites du système PoC. Proposer des pistes d'amélioration pour une version industrielle.}

\section{Perspectives d'industrialisation}
\todo{Décrire le chemin vers une version industrielle : intégration servo DS3218, robustesse, interface HMI, traçabilité, intégration au flux de production Pursuit.}

\section{Conclusion}
\todo{Synthèse des résultats et transition vers la conclusion générale.}
